\section{Introduction}

Symbolic regression (SR) searches over mathematical expressions, built from a small set of operators, variables and free constants, for the formula that best predicts an outcome $y$ from regressors $X$ \citep{schmidt2009,udrescu2020,cranmer2023}. Its promise for economics is an interpretable functional form found from the data rather than assumed: a demand curve, an Engel curve, a return to schooling, with elasticities that change with the level, thresholds, and diminishing returns read off the formula. The difficulty is that SR minimizes a prediction loss, so at best it is consistent for the conditional mean $\E[y \mid X]$. In the structural model
\[
y = f_0(X) + u, \qquad \E[u \mid X] \neq 0,
\]
the conditional mean is $f_0(X) + \E[u \mid X]$, and the formula SR returns describes the sum, not the law. In economics almost every regressor of interest is endogenous: prices are set in equilibrium with quantities, schooling is chosen, hours respond to wages, policies are taken up by those who benefit. SR applied to such data \citep{balla2025,rossetti2026} produces convincing formulas that are not causal.

Econometrics has three answers to endogeneity, and each loses something that SR offers. Parametric instrumental-variable (IV) estimation \citep{hansen1982} needs the functional form in advance. Sparse IV over a term library \citep{gautierrose,fanliao2014,linfengli2015,caner2009} selects terms under endogeneity, but only for laws that are linear in their parameters over a dictionary fixed in advance: $a + b\log x$ yes, $a x^{\theta}$ with unknown $\theta$, a kink at an unknown point, or a saturation curve $ax/(b + x)$ no. Nonparametric IV \citep{neweypowell2003,hallhorowitz2005,bck2007,chenpouzo2012} and its machine-learning versions \citep{hartford2017,singh2019,bennett2019,dikkala2020,xu2021} impose nothing on the form but return a curve or a network, not a formula, and the estimation problem is ill-posed. No method searches over free-form formulas while enforcing the conditional moment restriction $\E[y - f(X) \mid Z] = 0$ for an instrument $Z$. This paper fills that gap. What the moment restriction buys is that the formula found is a structural law and not a description of the conditional mean: the data choose the functional form, and the instrument carries the identification. The task is then to show that this can be done consistently, that it works at the sample sizes economists have, and that the method says so when the instruments do not identify the form.

The estimator proposed here has three steps. Candidate formulas are generated cheaply, by any means: a plain squared-error SR search, an SR search distilled from a sieve two-stage least-squares (2SLS) curve, a sparse path over a term library projected on the instruments, and a fixed library of textbook economic shapes with free constants. The constants of every candidate are then refitted by two-step GMM on $K$ sieve instruments, and the formula is selected by a GMM model-selection criterion in the sense of \citet{andrewslu2001}, Hansen's $J$ statistic plus a BIC penalty per constant, to which we add a per-node \emph{structure} penalty that accounts for the search over formula structures. Finally, an identification diagnostic built on the Anderson--Rubin statistic \citep{stockwright2000} reports the first-stage strength, identification-robust intervals for the constants, and the \emph{formula confidence set}: the candidate formulas the data do not reject. The verdict is ``identified'', ``weakly identified'' (with the formulas the data also allow listed) or ``rejected''.

The central finding of the Monte Carlo simulation is that the instrument does its work when formulas are \emph{scored}, not when they are generated. The squared-error search's own pick carries the endogeneity bias, but its list of candidate formulas usually contains the true form, and the moment criterion selects it. Distilling a sieve-2SLS curve gives candidates that select even better, and a fixed library of economic templates gives full coverage of the textbook shapes at no search cost. The expensive alternative, running the SR search itself on the moment criterion, buys nothing that scoring does not.

Section~\ref{sec:theory} gives the selection rule a theory. Bounded constants and bounded formula size turn the ill-posed inverse problem into a choice among finitely many parametric models, and the rule then selects a minimal true formula with probability tending to one for \emph{any} generator whose candidates contain one, with constants at the parametric rate and valid split-sample inference (\res{Theorem}{1}). A finite-sample bound (\res{Theorem}{2}) names the two ways selection fails, a wrong formula the moments barely separate from the truth and a longer formula that fits the moment noise, and shows that the per-node structure penalty is what controls the second as formulas grow. The formula confidence set covers the true form at the nominal level without any rank condition (\res{Proposition}{5}), so the diagnostic reports weak identification, rather than a confident formula, when the instruments do not identify the form. A $J$-point instrument cannot exclude a formula with $J$ or more constants whose moment equations have a solution within the constant bound (\res{Proposition}{4}), and under misspecification the estimator converges to the pseudo-true set, and to any Lipschitz law when the formula size grows with $n$ (\res{Theorem}{3}). The theory thus says what the Monte Carlo shows: once the candidates contain the truth, what decides is how far the moments separate them, not the search that produced them.

Section~\ref{sec:mc} reports a Monte Carlo of 1{,}098 SR searches over five designs (confounding, a heteroskedastic first stage that breaks the control-function approach, simultaneity, a discrete instrument, a weak instrument) and five laws (log, power, kink, saturation, two regressors), against nine baselines: oracle and wrong-form IV, sieve 2SLS, kernel IV, DeepIV, their symbolic distillations, library IV, control-function SR and an SRIT-style causal SR. With strong instruments the method recovers the log, power and two-regressor laws in 92--100\% of runs under confounding and a heteroskedastic first stage, and in 60--92\% under simultaneity and a discrete instrument; the squared-error search never recovers the power law under strong endogeneity, and where it finds the log form its slope is biased by exactly the amount the theory predicts. In the strong-endogeneity design recovery reaches 98--100\% from $n = 5000$ with the structure penalty, while with BIC alone recovery of the power law plateaus at 75--92\% because near-duplicates of the law win. Under the weak instrument recovery is 0--2\%, and the diagnostic flags every run. Section~\ref{sec:apps} applies the method to the Engel curves of the 1995 Family Expenditure Survey sample of \citet{bck2007}, where it reproduces the known shapes and reports no share as identified, and to \citet{angristkrueger1991}, where it reports weak identification instead of a formula, as it should. Section~\ref{sec:budget} shows that a fixed template library raises recovery where the moments separate the law, while an eight-fold search budget raises coverage but not recovery, because more candidates bring more wrong formulas that outscore the truth, the second channel of \res{Theorem}{2} in the data.

The closest related work is of three kinds, and the paper differs from each in one respect. Against sparse IV over a library, discussed above, the search here is over free-form formulas, and the library path becomes one candidate generator among others, whose limit turns out to be coverage rather than the IV logic: it never contains the power, kink or saturation laws. IV in system identification \citep{kuanglin2026,zeiringer2026} corrects coefficients for measurement noise over a fixed library, without external instruments and without a search over structures; here both are present. Causal constraints inside SR \citep{yan_srit,li2026,kacprzyk2024}, theory-constrained SR \citep{kronberger2022,cornelio2023,corywright2024,russeil2024} and sparse identification of dynamics \citep{brunton2016} restrict the inputs, the form or the library under causal sufficiency, ignorability or prior theory, and all still target $\E[y \mid X]$; unobserved confounding handled through instruments, which \citet{yan_srit} state as open, is what this paper adds. The nearest existing route, distilling an NPIV fit into a formula in the manner of \citet{cranmer2020}, is one of our generators, and the Monte Carlo locates its weakness in the selection step, which the criterion proposed here repairs. The econometric tools are standard: the criterion extends the GMM model selection of \citet{andrews1999} and \citet{andrewslu2001} to formula classes, the diagnostic applies \citet{stockwright2000} to a set of candidate formulas, the finite-sample bound uses \citet{tropp2012} and \citet{hsukakadezhang2012}, the misspecified and growing-size results follow the sieve literature \citep{aichen2003,chen2007,chenpouzo2012} and \citet{hallinoue2003}, and inference after selection is by sample splitting, with no uniform claim \citep{leebpotscher2005,berk2013}.

The assumptions are chosen to make the problem a selection problem, and each has a price. The error is additive and mean-independent of the instrument. This is what turns a candidate formula into a vector of moments that the data can score; it is also the form in which demand, Engel and earnings equations are written, and it excludes non-separable models with heterogeneous effects, for which the IV quantile model of \citet{chernozhukovhansen2005} is the natural extension. Constants are bounded and formulas have a maximal size. This turns the ill-posed inverse problem into a choice among finitely many parametric models and buys consistency and finite-sample bounds without a regularization scheme; its price is that rates of convergence for a growing size reduce to a modulus of injectivity on the formula sieve (\res{Proposition}{3.3}) that we do not bound, and bounding it for free-knot splines is the most tractable of the open problems. Two limits of the method itself are documented rather than fixed. A threshold law whose kink the instruments barely move and a backward bend at the edge of the data are almost never selected, because a simpler formula fits the moments equally well and no scoring rule can separate formulas the instruments do not distinguish. These are limits of identification within the formula class, not of the estimator; the diagnostic flags part of them, and the remedy is instruments that move the regressor through the kink or the bend, or far larger samples, not a larger search. Section~\ref{sec:limits} collects these with the remaining open problems.

The rest of the paper is organized as follows. Section~\ref{sec:model} states the model, the formula class and the estimator. Section~\ref{sec:theory} gives the theory, with proofs in Appendix~\ref{app:proofs}. Section~\ref{sec:impl} describes the implementation: the criteria, the candidate generators, the selection rules and the diagnostic. Sections~\ref{sec:mc}--\ref{sec:budget} report the Monte Carlo, the applications and the budget experiment. Section~\ref{sec:limits} concludes with the limits and open problems.
